
%\setcounter{page}{\value{page}}
\chapter{Introducere}
\label{cap:Introducere}
%\setcounter{page}{7}
%\addtocounter{page}{6}
\pagestyle{headings}
Ce se scrie aici:
\begin{itemize}
    \item Contextul
    \item Conturarea domeniului exact al temei
    \item Se răspunde la întrebările: \textbf{ce} (s-a făcut)?, \textbf{de ce} (s-a făcut, adică motivația; ce se rezolvă, la ce este util, etc.)?, \textbf{cum} (s-a făcut, adică particularitățile abordării, prezentate sumar).
    \item Introducerea se termină cu o descriere a conținutului lucrării, de genul: Cap X descrie ..., Cap Y prezintă ...
    \item Introducerea reprezintă o sinteză a lucrării, din care cititorul trebuie să-și poată da bine seama dacă lucrarea prezintă sau nu interes pentru el. 
\item reprezintă cca 5\% din lucrare
\end{itemize}

\section{Contextul proiectului}
Subcapitolul este scris cu font Times New Roman 14 puncte, bold. Pentru numerotare se folosesc cifre arabe, de ex: $X.Y$ unde X reprezintă numărul capitolului, iar Y numărul subcapitolului. Distanța de la titlul subcapitolului până la primul rând scris este de 8 puncte.

\subsection{Subsecțiune}
Distanța înainte și după titlul subcapitolului este de 8 puncte (acest lucru este realizat implicit în fișierul de stil furnizat). 
Fiecare tabel introdus în lucrare este numerotat astfel: Tabel $x.y$ unde x reprezintă numărul capitolului, iar y numărul tabelului din capitol. 
Se lasă o distanță de 8 puncte între text și tabel. (tabelul \ref{table:nonlin}).
\begin{table}[ht]
\centering                          % tabel centrat 
\begin{tabular}{|c|c|c|c|}          % 4 coloane centrate 
\hline\hline                        % linie orizontala dubla
Case & Method\#1 & Method\#2 & Method\#3 \\ [0.5ex]   % inserare tabel
%heading
\hline                              % linie orizontal simpla
1 & 50 & 837 & 970 \\               % corpul tabelului 
2 & 47 & 877 & 230 \\
3 & 31 & 25 & 415 \\[1ex]           % [1ex] adds vertical space
\hline                              
\end{tabular}
\caption{Nonlinear Model Results}   % titlul tabelului
\label{table:nonlin}                % \label{table:nonlin} introduce eticheta folosita pentru referirea tabelului in text; referirea in text se va face cu \ref{table:nonlin}
\end{table}

Fiecare figură introdusă în text este citată (de ex: în figura $x.y$ este prezentată ...) și numerotată.
Numerotarea se face automat astfel: Figura $x.y$ unde x reprezintă numărul capitolului, iar y numărul figurii din acel capitol. Ex: figura \ref{fig:imag}

\begin{figure}[h]
    \centering
    \includegraphics[]{image}	%numele fisierului
    \caption{Numele figurii}
    \label{fig:imag}
\end{figure}


Ecuații:
\begin{equation}
    \Delta =\sum_{i=1}^N w_i (x_i - \bar{x})^2 .
\end{equation}

\subsection{Alte exemple în \LaTeX}
\label{subsec:s10} %aceasta trebuie adaugat pentru a referi o sectiune  

Modus ponens modificat: 
% aceasta este o ecuatie 
\begin{equation} 
  \frac{a \land a \rightarrow b}
       {b}
\label{eq:e50}
\end{equation}

Acolada: 

\[\left \{ 
    \begin{array}{l}
        p =\frac{x}{x+y+z} \\
        q = \frac{y}{x+y+z}\\ 
        r = \frac{z}{x+y+z}
    \end{array}\right.
\]
        
Itemi ne-numerotați:  
\begin{itemize}
    \item numărul 1 
    \item numărul 2 
    \item numărul 3 
\end{itemize}

Itemi numerotați:
\begin{enumerate}
    \item numărul 1 
    \item numărul 2 
    \item numărul 3 
\end{enumerate}

%asa se citeaza o lucrare                    
Anumiți autori spun că  \cite{Zhou04} 


Un tabel mai complicat este \ref{Mom}.
\begin{table}
      \begin{center}
            \begin{tabular}{|l|l|l|l|l|l|l|}
                \hline 
                & \multicolumn{3}{c|}{belief} & \multicolumn{3}{c|}{atomicity} \\
                \cline{2-7}
                & $poor$ & $avg$ & $good$ & $poor$ & $avg$ & $good$ \\ 
        	\hline
                $success$   &  0.1    &  0.2    &  0.3    & 0.4    & 0.5    &  0.6     \\
                $failure$   &  0.1    &  0.2    &  0.3    & 0.4    & 0.5    &  0.6     \\
                                                    \hline
	    \end{tabular}
        \end{center}
    \caption{Multiplicarea Opiniilor Multinomiale}
    \label{Mom}
\end{table}

                                                                                                                                                                                                                                                                  
Un algoritm poate fi scris astfel:
%\usepackage{algorithm} si \usepackage{algorithmic}  trebuie incluse la inceputul documentului

\begin{algorithm}
 \caption{Commitment decision}
  \label{ftalg}
   \begin{algorithmic}
     \IF{$Committed(G_{1},GR,\alpha)$} 
      \STATE $BRT_{\alpha}$ = $PredictBRT(G_{1},GR,\alpha,C_{GR})$
      \STATE $C_{\beta}$ = $ContextUpdate(C_{\beta},o)$
      \STATE $BRT_{\beta}$ = $PredictBRT(G_{1},G_{2},\beta,C_{\beta})$
      \STATE $BRT_{\alpha}^{o}$ = $BRTReplace(BRT_{\alpha},BRT_{\beta})$
      \STATE $utility$ = $Eval(BRT_{\alpha}^{o})-Eval(BRT_{\alpha})$
     \ENDIF
      \IF{$utility \ge CommunicationCost(G_{2})$}
        \STATE $Int.To(G_{1},Communicate(G_{1},G_{2},o))$
      \ENDIF
   \end{algorithmic}
\end{algorithm}

\section*{Mulțumiri}
Mulțumiri lui Aristotel pentru Secțiunea \ref{subsec:s10} %Asa se refera o sectiune
